%!TEX TS-program = xelatex
%!TEX encoding = UTF-8 Unicode
%===============================================================================
%  MASTER CV  —  TECH / FAANG track
%  Target roles: ML Engineer, SWE-ML, Applied Scientist (ML), GenAI Engineer
%  READABILITY NOTES:
%   - Visible bullets = one line each, minimal bold (numbers + one marquee term).
%   - Deeper specifics live in "% INTERVIEW PREP" comments per entry.
%   - This is a *master* resume; trim/reorder per JD.
%===============================================================================
\documentclass[11pt, a4paper]{awesome-cv}
\geometry{left=1.4cm, top=1.35cm, right=1.4cm, bottom=1.35cm}
\usepackage{hyperref}

% -- Black body text + visible blue hyperlinks --
\definecolor{linkblue}{HTML}{0645AD}
\colorlet{text}{black}
\colorlet{graytext}{black}
\colorlet{darktext}{black}
\colorlet{lighttext}{black}
\colorlet{awesome}{black}
\hypersetup{colorlinks=true, urlcolor=linkblue, linkcolor=linkblue}

% -- Smaller, cleaner header --
\renewcommand*{\headerfirstnamestyle}[1]{{\fontsize{20pt}{1em}\headerfontlight\color{text} #1}}
\renewcommand*{\headerlastnamestyle}[1]{{\fontsize{20pt}{1em}\headerfont\bfseries\color{text} #1}}
\renewcommand*{\headerpositionstyle}[1]{{\fontsize{7.2pt}{1em}\bodyfont\scshape\color{text} #1}}
\renewcommand*{\headersocialstyle}[1]{{\fontsize{6.8pt}{1em}\bodyfont\color{linkblue} #1}}

% -- Tighter vertical spacing than stock, without cramping --
\renewcommand{\acvSectionTopSkip}{2.2mm}
\renewcommand{\acvSectionContentTopSkip}{3.0mm}
\renewcommand{\acvHeaderAfterSocialSkip}{4.0mm}
\setlength{\parskip}{6pt}

\name{Ishaan}{Gupta}
% \position{ML Engineer {\enspace}|{\enspace} Deep Learning Researcher}
\address{San Diego / Bay Area, CA}
\email{i3gupta@ucsd.edu}
\homepage{ishaanSD.github.io/home}
\linkedin{IgAI}
\github{IshaanSD}

\begin{document}
\makecvheader

%----------------------------------------
\begin{cvparagraph}
% CS PhD researcher in deep learning applications, building ML systems where success depends on more than the model: noisy measurements, limited labels, domain constraints, robust pipelines, and reliable evaluation.
Deep learning practitioner and researcher with domain expertise in healthcare, building robust AI pipelines for noisy, limited-label real-world data, with experience spanning biotechnology, hospital operations, recommendation systems, and production-grade software infrastructure.
\end{cvparagraph}

%----------------------------------------
\cvsection{Education}
\begin{cventries}
\cventry
{Ph.D. Computer Science — Deep Learning for Genomics \quad | \quad M.S. Computer Science, earned en route}
{University of California, San Diego}
{La Jolla, CA}
{Sep 2023 - Mar 2028 (exp.)}
{
\begin{cvitems}
\item {M.S. awarded Jun 2025 {\enspace}|{\enspace} GPA 3.92 {\enspace}|{\enspace} Charles Lee Powell Research Fellowship {\enspace}|{\enspace} Advisors: Tiffany Amariuta, Pavel A. Pevzner}
\end{cvitems}
}

\cventry
{B.S. Computer Science — Bioinformatics Specialization}
{University of California, San Diego}
{La Jolla, CA}
{Aug 2019 - Jun 2022}
{
\begin{cvitems}
\item {GPA 3.92 {\enspace}|{\enspace} Magna Cum Laude (Top 4\%) {\enspace}|{\enspace} IEEE Eta Kappa Nu {\enspace}|{\enspace} Minor in Biology}
\end{cvitems}
}
\end{cventries}

%----------------------------------------
\cvsection{Technical Skills}
\begin{cvskills}
  \cvskill
    {Languages}
    {Python, Java, C++, SQL, Bash, R}
  \cvskill
    {Deep Learning}
    {PyTorch, Lightning, Transformers, CNN, LSTM, PEFT / LoRA, FlashAttention, scikit-learn, XGBoost, LightGBM, TensorFlow}
  \cvskill
    {GenAI and Retrieval}
    {LangGraph, LangChain, MCP, RAG, tool-calling, ChromaDB, SentenceTransformers, agent evaluation}
  \cvskill
    {MLOps}
    {W\&B, MLflow, Hugging Face, Docker, FastAPI, REST APIs, CI/CD, pytest/JUnit, model evaluation, drift analysis}
  \cvskill
    {Data \& Infrastructure}
    {Spark, Kafka, Airflow, Pandas, Polars, SQL optimization, MongoDB, AWS, GCP, Slurm, Nextflow, Kubernetes}

\end{cvskills}


%----------------------------------------
\cvsection{Experience}
\begin{cventries}


  % =============================== DIPLOFORMER =================================
  % INTERVIEW PREP (not shown on CV):
  %  - 186M-param FlashAttention transformer; progressive unfreezing
  %    (backbone frozen epoch 1, then adapted).
  %  - Data: 400 individuals; EITHER 8,100 loci @4,096bp OR 2,700 loci @49kbp.
  %  - HW: mixed precision on 2×V100 or single H100; <100-sample runs -> 2 jobs
  %    co-located on one H100.
  %  - Nextflow DAG: CPU nodes (ingest/preprocess) -> GPU (train) -> light GPU (infer).
  %  - Fault tolerance: preempted Slurm jobs resume into the SAME W&B run.
  %  - Baselines Elastic Net + XGBoost. Drift detection stratified by ancestry and by
  %    biological features of genes/peaks (generalized to "subpopulations" on CV).
  %  - Predicts personalized molecular phenotypes (gene expression, chromatin
  %    accessibility) that generalize across loci AND across people.
  %  - Venues: Keystone 2025 (poster); ASHG 2026.
  \cventry
    {Graduate Student Researcher — Deep Learning | PyTorch, Nextflow, W\&B, MLflow, HuggingFace}
    {UCSD Computer Science Department}
    {San Diego, CA}
    {Sep 2024 - Present}
    {
      \begin{cvitems}
        \item {Built \textbf{DiploFormer}, a state-of-the-art framework for training and finetuning transformer models on population-level DNA sequence data.}
        \item {Fine-tuned \textbf{186M-parameter} FlashAttention \textbf{transformer} and CNN-based sequence models across individuals, loci, and molecular modalities.}
        \item {Stabilized pairwise training with output-accumulated losses and systematic sweeps of adaptation strategies and ablations.}
        \item {Orchestrated resource-aware CPU/GPU workflows with parallel jobs and checkpoint recovery across preprocessing, training, and inference.}
        \item {Benchmarked against Elastic Net/\textbf{XGBoost} baselines, monitored data drift, and validated disease-relevant regulatory signals.}
        \item {Finetuned LLM (Qwen 8B-parameter) using \textbf{QLoRA} with reasoning distillation from o4-mini, beating the teacher in accuracy (\textbf{61\% vs 48\%}) and latency (\textbf{100$\times$}) for gene perturbation and immune response prediction}
      \end{cvitems}
    }

  % =============================== INTUITIVE ===================================
  % INTERVIEW PREP (not shown on CV):
  %  - Work spans finetune I3D on kinetics (flow_imagenet.pt) - batches of 16, two inference modes - I3D classifier is good for annotations for which short-term memory suffices; I3D + LSTM is needed for annotations where long-term memory is useful.
  %  - Evaluated which annotations need 3D ConvNet alone vs which need LSTM; what temporal batch sizes are good (default 16). Research approches to inegrate 3d conv + lstm, https://github.com/google-deepmind/kinetics-i3d
  %  - LLM-assisted ingestion/validation, temporal outlier detection
  %  - LOF (Local Outlier Factor) ~O(n^2), memory-heavy. Replaced with GMM (linear)
  %    or boxplot/IQR (O(n log n)). Fill [N]× after benchmarking on one batch.
  %  - Latency = BATCH/compute (not endpoint): wall-clock of outlier stage before/
  %    after; full analytics-refresh time; memory drop. Serving is FastAPI but
  %    team-owned — mention only if role centers on serving.
  %  - "Trend algorithms + significance testing vs. downstream outcomes."
  %  - Product (DO NOT reveal): which OR / prep time / surgeon is slow; scheduling
  %    utilization. "Robust diagnostics" = median/IQR/MAD summaries.
  
  \cventry
    {PhD, AI/ML Data Scientist Intern | PyTorch, scikit-learn, SQL, LLM Analytics}
    {Intuitive Surgical, Inc.}
    {Sunnyvale, CA}
    {Jun 2026 - Present}
    {
      \begin{cvitems}
        \item {Expanding the end-to-end AI analytics pipeline from 3D ConvNet and LSTM-based action-recognition to SQL-based metrics generation and LLM-powered customer-facing insights.}
        \item {Optimizing model architecture choice, temporal batches, hyperparameters, and inference platform for classification on camera streams.}
        \item {Implemented vector quantization and temporal outlier detection for robust and interpretable workflow recommendation scores.}
        \item {Built an LLM-assisted ingestion and validation workflow to standardize hospital observation spreadsheets for user research and trend analysis.}
      \end{cvitems}
    }

  % =============================== ABTERRA ====================================
  % INTERVIEW PREP: 4× = improvement on UNSEEN, company-specific data after
  % fine-tuning autoregressive transformers for de novo peptide sequencing (mass-spec).
  \cventry
    {Machine Learning Engineer Intern | PyTorch Lightning, W\&B, Java/DL4J}
    {Abterra Biosciences}
    {San Diego, CA}
    {Jul 2024 - Sep 2024}
    {
      \begin{cvitems}
        \item {Finetuned autoregressive transformer models for de novo peptide sequencing, improving held-out performance \textbf{4$\times$} on in-house data.}
        \item {Built reproducible MLOps workflows for hyperparameter sweeps, model comparison, failure-mode analysis, and sequence bias detection.}
        \item {Achieved \textbf{95\% accuracy} (baseline: 60\%) in resolving ambiguous antibodies through peptide–spectrum feature engineering and FNN classifiers.}
        \item {Deployed trained deep learning models into the company's Java codebase for parallelized performance evaluation.}
      \end{cvitems}
    }
    
  % =============================== PEVZNER ====================================
  % INTERVIEW PREP:
  %  - Repeat-finding project: approximate + suffix-array pattern matching with an
  %    ITERATIVE de Bruijn graph simplification procedure; repeat-motif lengths
  %    10,000 to 1,000,000 characters. Beat SOTA repeat detection in the critical
  %    immune-system DNA loci (Primate T2T; Nature 2025). Suffix-array alignment
  %    (unialigner) optimized 4-10× in C++.
  %  - Signal-processing (SEPARATE project): cumulative sums (CUSUM) for nanopore
  %    blockade detection; FFT + clustering to infer directionality of the peptide
  %    sequence.
  %  - Papers: Nature (ape genomes / Primate T2T) + Bioinformatics (GenomeDecoder).
  %    Plotly genome-scale viz lives on the Bio CV.
  \cventry
    {Graduate Student Researcher — Algorithms | C++, Graph Algorithms, Suffix-array, Signal Processing}
    {UCSD Computer Science Department}
    {San Diego, CA}
    {Jan 2022 - Jul 2024}
    {
      \begin{cvitems}
        \item {Outperformed state-of-the-art algorithms for segmental duplication detection; built scalable visualization tool. \textit{Bioinformatics} (2025).}
        \item {Co-designed GenomeDecoder, that iteratively simplifies de Bruijn graphs to find 10K-1M long inexact repeats in 100M-base DNA sequences.}
        \item {Improved runtime of kmer-count-stratified sequence alignment approach (UniAligner) by \textbf{4-10$\times$} and integrated it with GenomeDecoder to discover gene deletions across primate genomes. \textit{Nature} (2025).}
        \item {Developed \textbf{signal-processing} methods using cumulative sums, FFT, and clustering for bidirectionality detection in noisy Nanopore signals.}
      \end{cvitems}
    }
    
  % =============================== ILLUMINA ===================================
  % INTERVIEW PREP: Agile; LIMS = Lab Information Management System; genotyping
  % 10,000s samples/day; Java Spring + SQL + Angular; REST, AWS, RabbitMQ; containerized.
  \cventry
    {Software Engineer I — Systems Integration | Java Spring, AWS, Docker, Kubernetes, RabbitMQ, TypeScript}
    {Illumina, Inc.}
    {San Diego, CA}
    {Jun 2022 - Dec 2022}
    {
      \begin{cvitems}
        \item {Developed production-grade ETL infrastructure code for high-throughput lab automation workflows processing \textbf{10{,}000s of samples/day}.}
        \item {Built Java Spring/JDBC REST APIs and event-driven system integrating containerized microservices, and contributed to user-facing UI.}
        \item {Deployed \textbf{Dockerized} services on Kubernetes with SQL stored procedures, JUnit tests, and CI/CD release workflows.}
      \end{cvitems}
    }

  \cventry
    {ML Project Consultant | Keras, scikit-learn, BigQuery, Beautiful Soup (bs4), spaCy}
    {Model Medicines}
    {San Diego, CA}
    {Nov 2020 - Mar 2021}
    {
      \begin{cvitems}
        \item {Engineered FNN+Random Forest ensemble model and \textbf{GCP BigQuery}-powered SQL pipelines for COVID-19 drug-discovery screens.}
        \item {Led a team building an NLP triage workflow that mines drug candidates and their synonyms in the scientific literature, and predicts novelty.}
      \end{cvitems}
    }

\end{cventries}

%----------------------------------------
\cvsection{Selected Projects}
\begin{cventries}
% GRAPHAZON — recommendation / ranking. Learn-to-rank via LambdaMART (LightGBM);
    % NDCG/MRR; Amazon ESCI Shopping Queries; KG + Node2Vec + text embeddings.
    \cventry
    {LangChain · Knowledge Graphs · LightGBM / LambdaMART}
    {Graphazon - Natural Language Query-based Product Recommendation}
    {~}
    {(\href{https://github.com/IshaanSD/Graphazon}{github.com/IshaanSD/Graphazon})}
    {
      \begin{cvitems}
        \item {Built a product-recommendation pipeline on the Amazon ESCI dataset, using an LLM to parse queries into typed intents.}
        \item {Combined knowledge-graph, Node2Vec, and text embeddings with \textbf{LambdaMART} learn-to-rank optimized for Mean Reciprocal Rank.}
      \end{cvitems}
    }

% MITOEDIT — MCP server lets any agent run tool via NL; multi-format ingestion;
  % downstream queries e.g. "targets with no non-benign bystanders". Existing stack:
  % Python + FastAPI web UI + Docker + tests + alignment/liftover -> CSVs. bioRxiv preprint.
  \cventry
    {MCP · LLM Agents · FastAPI · Docker · Python}
    {MitoEdit — Lab-in-the-Loop Agent for Mitochondrial Research at St. Jude}
    {~}
    {(\href{https://github.com/xz-stjude/MitoEdit}{github.com/xz-stjude/MitoEdit})}
    {
        \begin{cvitems}
            \item {Identifies mtDNA cancer variants editable by TALE and predicts bystander effects through user-defined rules. \textit{(manuscript in preparation)}}
            \item {Added an \textbf{MCP server} and FastAPI layer so bench scientists can use by prompting or UI. Built a dockerized, CI/CD-enabled Python workflow.}
        \end{cvitems}
    }
    
  % CHEMBOT — 3 MCP servers over openFDA (labels + FAERS) and PubMed; LangGraph
  % multi-route router; ChromaDB + SentenceTransformers + OpenAI; deployed on GCP.
    \cventry
    {LangGraph · MCP · RAG · ChromaDB · Google Cloud}
    {{ChemBot - Agentic RAG Assistant for Trusted OTC Drug Information}}
    {~}
    {(\href{https://github.com/IshaanSD/ChemBot}{github.com/IshaanSD/ChemBot})}
    {
        \begin{cvitems}
            \item {Agentic system for over-the-counter drug information retrieval (eg: side-effects) grounded in trusted biomedical evidence (FAERS and PubMed).}
            \item {\textbf{LangGraph} routing across MCP servers integrated with RAG (\textbf{ChromaDB}), embeddings (\textbf{SentenceTransformer}), and LLMs (OpenAI/Gemini)}
            \item {Deployed on Google Cloud with automated evaluations for tool-calling, agent skills, citation faithfulness, and cross-LLM response quality.}
        \end{cvitems}
    }
    
  % GENERATIVE MODELS — latent diffusion (UNet, conditional chest X-ray), β-VAE
  % (interpretable latents), UNet-GAN (mammogram augmentation + tumor classification).
  \cventry
    {Diffusion Model · $\beta$-VAE · Context U-Net · GAN}
    {Generative Models for Medical Image Augmentation}
    {~}
    {(\href{https://github.com/IshaanSD/DLCXR}{github.com/IshaanSD/DLCXR})}
    {
        \begin{cvitems}
            \item {Built a \textbf{diffusion model} training pipeline using $\beta$-VAE latent encoding and a Context U-Net denoiser, and demonstrated on chest X-ray images.}
            \item {Implemented disentangled $\beta$-VAE representations for CXR generation and a U-Net GAN baseline for mammography augmentation.}
        \end{cvitems}
    }
\end{cventries}

%----------------------------------------
\cvsection{Publications \& Conferences}
\begin{cventries}
\cventry
{~}{~}{~}{~}
{
\begin{cvitems}
% \item {\textbf{Gupta I.}, Amariuta T. \textit{DiploFormer: Deep Learning framework for capturing Non-Additive Regulatory Effects on Gene Expression and Chromatin Accessibility.} Presented at Keystone AI in Molecular Biology, 2025; \textbf{ASHG}, 2026. (\href{https://github.com/IshaanSD/DiploFormer_1000G_atac}{github.com/IshaanSD/DiploFormer\_1000G\_atac})}
% \item {Zhang Z., \textbf{Gupta I.}, Pevzner P.A. \textit{GenomeDecoder: inferring segmental duplications in repetitive genomic regions.} \textbf{Bioinformatics}, 2025.}
% \item {Yoo D.A., et al. \textit{Complete sequencing of ape genomes.} \textbf{Nature}, 2025.} (\href{https://ishaansd.github.io/iPlots}{ishaansd.github.io/iPlots})
\item {\textit{DiploFormer: Deep Learning framework for capturing Non-Additive Regulatory Effects on Gene Expression and Chromatin Accessibility.} Presented at Keystone AI in Molecular Biology, 2025; \textbf{ASHG}, 2026. (\href{https://github.com/IshaanSD/DiploFormer_1000G_atac}{github.com/IshaanSD/DiploFormer\_1000G\_atac})}
\item {\textit{GenomeDecoder: inferring segmental duplications in repetitive genomic regions. (\href{https://ishaansd.github.io/iPlots}{https://github.com/IshaanSD/bestFitAl})
\textbf{Bioinformatics}, 2025.}}
\item {\textit{Complete sequencing of ape genomes.} \textbf{Nature}, 2025. (\href{https://ishaansd.github.io/iPlots}{ishaansd.github.io/iPlots})}
\end{cvitems}
}
\end{cventries}

\cvsection{Certifications}
\begin{cventries}
\cventry
{~}{~}{~}{~}
{
\begin{cvitems}
\item {\textbf{AI Agents:} Google AI Agents Capstone \quad | \quad \textbf{Data Engineering:} IBM SQL, ETL/Airflow/Kafka, Spark/Hadoop \quad | \quad \\ \textbf{Backend:} Advanced Java Spring Integration Testing}
\end{cvitems}
}
\end{cventries}

\cvsection{Leadership \& Mentoring}
\begin{cventries}
  \cventry
    {Vice President and Workshop Chair – Undergraduate Bioinformatics Club}
    {UC San Diego}
    {}
    {Aug 2020 – Jul 2022}
    {
      \begin{cvitems}
        \item {Led 15-member team hosting research talks, mentorship workshops, and social events. Mentored students on learning to code and do ML.}
        \item {Developed 11 Bioinformatics / Machine Learning workshops and deployed on AWS EC2 to make students capable of contributing to computational research.
        (\href{bioinformaticscrashcourse.com}{bioinformaticscrashcourse.com})}
      \end{cvitems}
    }

\cventry
{Teaching Assistant / Tutor}
{UC San Diego}
{San Diego, CA}
{Sep 2020 - Jun 2022}
{
\begin{cvitems}
\item {Courses: \textit{CSE 181}, \textit{CSE 182}, \textit{BICD 140}, \textit{BILD 4}; taught algorithms, machine learning concepts, statistics for research, and biological data analysis.}
\item {Supported instruction on VAEs, Transformers, HMMs, clustering, graph algorithms, dynamic programming, databases, and hypothesis testing.}
\end{cvitems}
}
\end{cventries}

% \cvsection{Teaching Experience}
% \begin{cventries}
% \cventry
% {Teaching Assistant / Tutor}
% {UC San Diego}
% {San Diego, CA}
% {Sep 2020 - Jun 2022}
% {
% \begin{cvitems}
% \item {Taught data structures and algorithms, machine learning, biological databases, immunology, and lab methods across 4 courses, supporting technical instruction, grading, and student mentoring.}
% \item {Explained complex topics including graph algorithms, dynamic programming, HMMs, clustering, statistics, and scientific writing to undergraduate students.}
% \end{cvitems}
% }
% \end{cventries}


\end{document}
